\section{Phase 1 : Optimisation des caractéristiques}
\label{sec:phase1}

\subsection{Stratégie}

Le LBP, la DWT et le HOG exposent chacun des hyperparamètres ajustables. Entraîner
un classifieur complet pour chaque combinaison nécessiterait des milliers de modèles.
À la place, nous classons les configurations de paramètres à l'aide de trois
\emph{métriques de séparabilité} calculées directement sur les vecteurs de caractéristiques,
sans entraînement de classifieur :

\begin{itemize}
  \item \textbf{Ratio de Fisher (FDR) :} rapport de la variance inter-classe sur la variance intra-classe.
  \item \textbf{Information mutuelle (MI) :} dépendance statistique entre caractéristiques et étiquettes de classe.
  \item \textbf{Indice de Davies--Bouldin (DBI) :} chevauchement des clusters dans l'espace des caractéristiques
        (plus bas est meilleur ; nous l'inversons avant agrégation).
\end{itemize}

Un score composite est obtenu par normalisation min--max de chaque métrique sur la grille,
puis moyennage :
\[
  \text{Score} = \frac{1}{3}\Big(\widehat{\text{FDR}} + \widehat{\text{MI}} + (1 - \widehat{\text{DBI}})\Big)
\]

La configuration atteignant le score composite le plus élevé par famille de descripteur
est figée pour toutes les expériences de Phase~2. L'espace de recherche total couvre
\NPhaseOneConfigs{} configurations.

\subsection{Optimisation du LBP}

La grille de recherche couvrait les voisins $P \in \{4,8,12,16\}$, le rayon $R \in \{1,2,3\}$
et trois méthodes d'encodage (\texttt{uniform}, \texttt{nri\_uniform}, \texttt{ror}).
La Figure~\ref{fig:lbp_search} présente le paysage de séparabilité.

\begin{figure}[H]
  \centering
  \includegraphics[width=0.72\linewidth]{lbp_param_search.png}
  \caption{Recherche de paramètres LBP : score de séparabilité selon la configuration.
           La meilleure configuration (P=12, R=1, uniform) atteint un score de \lbpScore.}
  \label{fig:lbp_search}
\end{figure}

Parmi les méthodes d'encodage, \texttt{uniform} surpasse systématiquement les alternatives
pour tous les couples $(P,R)$ (Figure~\ref{fig:lbp_method}). L'histogramme compact
des motifs à transitions limitées réduit le bruit et améliore la généralisation
entre patients.

\begin{figure}[H]
  \centering
  \includegraphics[width=0.72\linewidth]{lbp_by_method.png}
  \caption{Comparaison des méthodes d'encodage LBP. L'encodage \texttt{uniform} obtient
           la meilleure séparabilité pour la plupart des combinaisons de paramètres.}
  \label{fig:lbp_method}
\end{figure}

\subsection{Optimisation de la DWT}

Sept familles d'ondelettes (haar, db1, db2, db4, sym2, coif1, bior1.3) aux niveaux
$\{1,2,3\}$ ont été évaluées (21 configurations). La Figure~\ref{fig:dwt_search}
présente les résultats.

\begin{figure}[H]
  \centering
  \includegraphics[width=0.72\linewidth]{dwt_param_search.png}
  \caption{Recherche de paramètres DWT selon les familles d'ondelettes et les niveaux
           de décomposition. \texttt{db2} au niveau~1 atteint la meilleure séparabilité (\dwtScore).}
  \label{fig:dwt_search}
\end{figure}

Une comparaison directe des familles d'ondelettes (Figure~\ref{fig:dwt_wavelets}) confirme
que \texttt{db2} offre le meilleur compromis : elle capture les motifs tissulaires
tumoraux en moyenne fréquence sans sur-fragmenter le signal comme le font les familles
plus complexes.

\begin{figure}[H]
  \centering
  \includegraphics[width=0.72\linewidth]{dwt_wavelet_comparison.png}
  \caption{Comparaison des familles d'ondelettes au niveau de décomposition optimal.
           \texttt{db2} surpasse les familles plus simples (haar) et plus complexes (bior1.3).}
  \label{fig:dwt_wavelets}
\end{figure}

\subsection{Optimisation du HOG}

Le HOG a été exploré sur 12 configurations variant le nombre d'orientations (6/9),
la taille des cellules (8/16 pixels) et la taille des blocs (2/3 cellules). Une
configuration avec 6 orientations, cellules $(16\times16)$, blocs $(3\times3)$ et
normalisation L2-Hys a atteint la séparabilité la plus élevée (\hogScore).

\subsection{Synthèse de la Phase 1}

Le Tableau~\ref{tab:phase1_summary} rapporte les paramètres retenus et leurs scores
de séparabilité. La Figure~\ref{fig:phase1_summary} fournit une synthèse visuelle.

\begin{table}[H]
  \centering
  \caption{Paramètres retenus en Phase~1 et scores de séparabilité.}
  \label{tab:phase1_summary}
  \begin{tabular}{@{}llcl@{}}
    \toprule
    \textbf{Descripteur} & \textbf{Meilleurs paramètres} & \textbf{Score} & \textbf{Branche} \\
    \midrule
    Statistiques & \statsParams & --- & A \\
    LBP & \lbpParams & \lbpScore & B \\
    DWT & \dwtParams & \dwtScore & C \\
    HOG & \hogParams & \hogScore & D \\
    \bottomrule
  \end{tabular}
\end{table}

\begin{figure}[H]
  \centering
  \includegraphics[width=0.80\linewidth]{phase1_summary.png}
  \caption{Synthèse de l'optimisation Phase~1 pour les quatre familles de descripteurs.}
  \label{fig:phase1_summary}
\end{figure}

Les descripteurs de texture (LBP) et de gradient (HOG) atteignent des scores de
séparabilité quasi parfaits, ce qui en fait les branches individuelles les plus
fortes. Les statistiques offrent une base légère ; la DWT apporte une information
fréquentielle complémentaire.
